% Bagian: first-order-logic
% Bab: syntax-and-semantics
% Subbagian: unique-readability

\documentclass[../../../include/open-logic-section]{subfiles}

\begin{document}

\olfileid[id]{fol}{syn}{unq}

\olsection{Keterbacaan Unik}

\begin{explain}
Cara kita mendefinisikan !!{formula}s menjamin bahwa setiap !!{formula}
memiliki \emph{pembacaan unik}, yakni pada dasarnya hanya ada satu cara
untuk membangunnya menurut aturan pembentukan !!{formula}s dan hanya satu
cara untuk ``menafsirkannya.'' Jika tidak demikian, kita akan mempunyai
!!{formula}s ambigu, yakni !!{formula}s yang memiliki lebih dari satu
pembacaan atau penafsiran---dan hal ini jelas ingin kita hindari. Lebih
penting lagi, tanpa sifat ini, sebagian besar definisi dan bukti yang akan
kita berikan tidak akan berhasil.

Barangkali cara terbaik untuk memperjelas hal ini ialah melihat apa yang
akan terjadi jika kita memberikan aturan pembentukan !!{formula}s yang buruk
dan tidak menjamin keterbacaan unik. Misalnya, kita dapat saja melupakan
tanda kurung dalam aturan pembentukan untuk penghubung; sebagai contoh,
kita mungkin mengizinkan aturan berikut:
\begin{quote}
Jika $!A$ dan $!B$ merupakan !!{formula}s, maka $!A \lif !B$ juga
merupakan !!a{formula}.
\end{quote}
Berawal dari formula atomik $!D$, aturan ini memungkinkan kita membentuk
$!D \lif !D$. Dari formula tersebut bersama $!D$, kita akan memperoleh
$!D \lif !D \lif !D$. Namun, ada dua cara untuk melakukannya:
\begin{enumerate}
\item Kita mengambil $!D$ sebagai $!A$ dan $!D \lif !D$ sebagai $!B$.
\item Kita mengambil $!A$ sebagai $!D \lif !D$ dan $!B$ sebagai~$!D$.
\end{enumerate}
Sejalan dengan itu, ada dua cara untuk ``membaca'' !!{formula}~$!D \lif !D
\lif !D$. Formula tersebut berbentuk $!B \lif !C$ dengan $!B$ ialah $!D$
dan $!C$ ialah $!D \lif !D$, tetapi formula itu \emph{juga} berbentuk $!B
\lif !C$ dengan $!B$ ialah $!D \lif !D$ dan $!C$ ialah~$!D$.

Jika hal ini terjadi, definisi kita tidak selalu berhasil. Misalnya, ketika
kita mendefinisikan !!{main operator} suatu formula, kita mengatakan: dalam
formula yang berbentuk $!B \lif !C$, !!{main operator} adalah kemunculan
$\lif$ yang ditunjukkan. Namun, jika formula $!D \lif !D \lif !D$ dapat
dicocokkan dengan $!B \lif !C$ melalui dua cara berbeda yang disebutkan di
atas, dalam kasus pertama kita memperoleh kemunculan pertama $\lif$ sebagai
!!{main operator}, sedangkan dalam kasus kedua kita memperoleh kemunculan
kedua. Padahal, kita bermaksud menjadikan !!{main operator} sebagai suatu
\emph{fungsi} dari !!{formula}; dengan kata lain, setiap !!{formula} harus
memiliki tepat satu kemunculan !!{main operator}.
\end{explain}

\begin{lem}
Banyaknya tanda kurung kiri dan tanda kurung kanan dalam
!!a{formula}~$!A$ adalah sama.
\end{lem}

\begin{proof}
Kita membuktikannya melalui induksi pada cara $!A$ dibangun. Hal ini
memerlukan dua langkah: (a) Mula-mula kita harus membuktikan bahwa semua
formula atomik memiliki sifat yang dimaksud (basis induksi). (b) Kemudian
kita harus membuktikan bahwa ketika kita membangun formula baru dari formula
yang diberikan, formula baru itu memiliki sifat tersebut asalkan formula
lama memilikinya.

Misalkan $l(!A)$ menyatakan banyaknya tanda kurung kiri dan $r(!A)$
banyaknya tanda kurung kanan dalam~$!A$; demikian pula, $l(t)$ dan $r(t)$
masing-masing menyatakan banyaknya tanda kurung kiri dan kanan dalam suatu
suku~$t$.

\begin{prob}
    Buktikan bahwa untuk setiap suku~$t$, $l(t) = r(t)$.
\end{prob}

\begin{enumerate}
\tagitem{prvFalse}{\indcase{!A}{\lfalse}{$\indfrm$ memiliki $0$ tanda
    kurung kiri dan $0$ tanda kurung kanan.}}{}

\tagitem{prvTrue}{\indcase{!A}{\ltrue}{$\indfrm$ memiliki $0$ tanda
    kurung kiri dan $0$ tanda kurung kanan.}}{}

\item \indcase{!A}{\Atom{R}{t_1,\dots,t_n}}{$l(\indfrm) = 1 + l(t_1) +
  \dots + l(t_n) = 1 + r(t_1) + \dots + r(t_n) = r(\indfrm)$. Di sini kita
  menggunakan fakta, yang diserahkan sebagai latihan, bahwa $l(t) = r(t)$
  untuk setiap suku~$t$.}

\item \indcase{!A}{\eq[t_1][t_2]}{$l(\indfrm) = l(t_1) + l(t_2) =
  r(t_1) + r(t_2) = r(\indfrm)$.}

\tagitem{prvNot}{\indcase{!A}{\lnot !B}{Menurut hipotesis induksi,
    $l(!B) = r(!B)$. Jadi, $l(\indfrm) = l(!B) = r(!B) =
    r(\indfrm)$.}}{}

\item \indcase{!A}{(!B \ast !C)}{Menurut hipotesis induksi, $l(!B) =
  r(!B)$ dan $l(!C) = r(!C)$. Jadi, $l(\indfrm) = 1 + l(!B) + l(!C) =
  1 + r(!B) + r(!C) = r(\indfrm)$.}

\tagitem{prvAll}{\indcase{!A}{\lforall[x][!B]}{Menurut hipotesis
    induksi, $l(!B) = r(!B)$. Jadi, $l(\indfrm) = l(!B) = r(!B) =
    r(\indfrm)$.}}{}

% Cetak kasus \lexists jika prvEx dan bukan prvAll
\tagitem{prvAll,notprvEx}{}{\indcase{!A}{\lexists[x][!B]}{Menurut hipotesis
    induksi, $l(!B) = r(!B)$. Jadi, $l(\indfrm) = l(!B) = r(!B) =
    r(\indfrm)$.}}

% Cukup katakan serupa jika prvEx dan prvAll
\tagitem{notprvAll,notprvEx}{}{\indcase{!A}{\lexists[x][!B]}{Serupa.}}
\end{enumerate}
\end{proof}

\begin{defn}[Prefiks sejati]
Suatu untai simbol $!B$ merupakan \emph{prefiks sejati} dari suatu untai
simbol~$!A$ jika konkatenasi $!B$ dengan suatu untai simbol takkosong
menghasilkan~$!A$.
\end{defn}

\begin{lem}\ollabel{lem:no-prefix}
Jika $!A$ merupakan !!a{formula} dan $!B$ merupakan prefiks sejati dari
$!A$, maka $!B$ bukan !!a{formula}.
\end{lem}

\begin{proof}
Latihan.
\end{proof}

\begin{prob}
Buktikan \olref[fol][syn][unq]{lem:no-prefix}.
\end{prob}

\begin{prop}
\ollabel{prop:unique-atomic}
Jika $!A$ merupakan !!{formula} atomik, maka formula tersebut memenuhi tepat
satu dari kondisi-kondisi berikut.
\begin{enumerate}
\tagitem{prvFalse}{$!A \ident \lfalse$.}{}
\tagitem{prvTrue}{$!A \ident \ltrue$.}{}
\item $!A \ident \Atom{R}{t_1,\dots,t_n}$, dengan $R$ merupakan
  !!{predicate} $n$-tempat, $t_1$, \dots, $t_n$ merupakan suku, dan
  masing-masing dari $R$, $t_1$, \dots, $t_n$ ditentukan secara unik.
\item $!A \ident \eq[t_1][t_2]$, dengan $t_1$ dan $t_2$ merupakan suku
  yang ditentukan secara unik.
\end{enumerate}
\end{prop}

\begin{proof}
Latihan.
\end{proof}

\begin{prob}
Buktikan \olref[fol][syn][unq]{prop:unique-atomic} (Petunjuk: Rumuskan dan
buktikan versi \olref[fol][syn][unq]{lem:no-prefix} untuk suku.)
\end{prob}

\begin{prop}[Keterbacaan Unik]
Setiap !!{formula} memenuhi tepat satu dari kondisi-kondisi berikut.
\begin{enumerate}
\item $!A$ bersifat atomik.

\tagitem{prvNot}{$!A$ berbentuk $\lnot !B$.}{}

\tagitem{prvAnd}{$!A$ berbentuk $(!B \land !C)$.}{}

\tagitem{prvOr}{$!A$ berbentuk $(!B \lor !C)$.}{}

\tagitem{prvIf}{$!A$ berbentuk $(!B \lif !C)$.}{}

\tagitem{prvIff}{$!A$ berbentuk $(!B \liff !C)$.}{}

\tagitem{prvAll}{$!A$ berbentuk $\lforall[x][!B]$.}{}

\tagitem{prvEx}{$!A$ berbentuk $\lexists[x][!B]$.}{}
\end{enumerate}
Selain itu, dalam setiap kasus, $!B$, atau $!B$ dan $!C$, ditentukan secara
unik. Artinya, misalnya, tidak ada pasangan berbeda $!B$, $!C$ dan $!B'$,
$!C'$ sedemikian sehingga $!A$ sekaligus berbentuk
\iftag{prvIf}{$(!B \lif !C)$ dan $(!B' \lif !C')$.}{
    \iftag{prvOr}{$(!B \lor !C)$ dan $(!B' \lor !C')$.}{
      \iftag{prvAnd}{$(!B \land !C)$ dan $(!B' \land !C')$.}{}}}
\end{prop}

\begin{proof}
Aturan pembentukan mensyaratkan bahwa jika !!a{formula} tidak atomik,
formula itu harus diawali oleh tanda kurung buka~(,
\iftag{prvNot}{$\lnot$,}{} atau suatu kuantor. Di sisi lain, setiap
!!{formula} yang diawali oleh salah satu simbol berikut harus bersifat
atomik: !!a{predicate}, !!a{function},
!!a{constant}\iftag{prvFalse}{, $\lfalse$}{}\iftag{prvTrue}{, $\ltrue$}{}.

Jadi, kita hanya perlu menunjukkan bahwa jika $!A$ berbentuk $(!B \ast
!C)$ dan juga berbentuk $(!B' \mathbin{\ast'} !C')$, maka $!B \ident !B'$,
$!C \ident !C'$, dan $\ast = {\ast'}$.

Andaikan $!A \ident (!B \ast !C)$ dan $!A \ident (!B'
\mathbin{\ast'} !C')$ keduanya berlaku. Maka $!B \ident !B'$ atau tidak.
Jika berlaku, jelaslah bahwa $\ast = {\ast'}$ dan $!C \ident !C'$, karena
keduanya kemudian merupakan subuntai $!A$ yang dimulai pada posisi yang sama
dan memiliki panjang yang sama. Kasus lainnya ialah $!B \not\ident !B'$.
Karena $!B$ dan $!B'$ sama-sama merupakan subuntai $!A$ yang dimulai pada
posisi yang sama, salah satunya harus menjadi prefiks sejati dari yang lain.
Namun, hal ini mustahil menurut \olref{lem:no-prefix}.
\end{proof}


\end{document}
